\chapter{Complex numbers: amplitude and phase}\label{ch:complex}
\begin{readergoal}Build the complex plane from real coordinates and understand Euler's formula before meeting Fourier and Hilbert transforms.\end{readergoal}
\section{Pairs with new arithmetic}
A complex number $z=a+ib$ has real part $a$ and imaginary part $b$, where $i^2=-1$. We may regard $z$ as the ordered pair $(a,b)$ of reals, with addition $(a,b)+(c,d)=(a+c,b+d)$ and multiplication $(a,b)(c,d)=(ac-bd,ad+bc)$. The conjugate is $\bar z=a-ib$ and $|z|=\sqrt{z\bar z}=\sqrt{a^2+b^2}$. For $z\ne0$, $z^{-1}=\bar z/|z|^2$.

\section{A plane with angles}
Write $z=r(\cos\theta+i\sin\theta)$ for modulus $r\ge0$ and argument $\theta$. Angles differing by integer multiples of $2\pi$ represent the same direction. Euler's formula $e^{i\theta}=\cos\theta+i\sin\theta$ gives the polar representation $z=re^{i\theta}$. Multiplication adds arguments and multiplies moduli:
\[(r_1e^{i\theta_1})(r_2e^{i\theta_2})=r_1r_2e^{i(\theta_1+\theta_2)}.\]
The periodicity of phase will require circular, rather than ordinary linear, statistics.

\section{Why complex numbers are indispensable for signals}
A sine wave can be written $x(t)=A\cos(2\pi ft+\phi)=\Re\{Ae^{i(2\pi ft+\phi)}\}$. Multiplication by $e^{i\omega t}$ shifts the angular coordinate at a constant rate. The Fourier transform treats a signal as a superposition of such rotating complex exponentials. The complex coefficient encodes two quantities: magnitude and phase. Discarding one generally loses information about the waveform.

\begin{workedexample}[Phase wraps]
The angles $179^\circ$ and $-179^\circ$ differ numerically by $358^\circ$ but lie just $2^\circ$ apart on a circle. Their arithmetic mean is $0^\circ$, which points in the opposite direction. A circular mean instead uses the argument of $e^{i179^\circ}+e^{-i179^\circ}$ and correctly points near $180^\circ$.
\end{workedexample}
\begin{exercises}
\item Compute $(1+i)^2$ and $|1+i|$.
\item Show $|zw|=|z||w|$ using complex conjugation.
\item Express $\cos\theta\cos\phi$ as a combination of complex exponentials and simplify.
\item Explain why phase averages need not be meaningful when oscillatory amplitude is close to zero.
\end{exercises}